% =====================================================================
% ANEXO A -- MODELO DE CARTEL INFORMATIVO
% Responsable: Fabian Moreno Ugarte
%
% El cartel se maqueta con \fbox + minipage + tabular, sin tcolorbox ni
% mdframed, para que compile en cualquier distribucion sin instalar
% paquetes adicionales. Desde 2026-09-14 va en negro, como todo el
% informe: el marco se conserva porque representa el borde fisico del
% cartel, no es un recuadro de aviso.
% =====================================================================
\section{Modelo de cartel informativo de videovigilancia}
\label{anx:cartel}

\subsection{Propósito y estatus de este modelo}

La Sección~\ref{sec:privacidad-diseno} estableció que el tratamiento de datos
personales mediante videovigilancia exige informar previamente a los
titulares mediante carteles visibles. Este anexo propone un modelo
concreto para el sistema del proyecto.

Se trata de una \textbf{propuesta de
trabajo}, no de un formato aprobado. Está construido sobre los elementos
informativos que la Directiva 01-2020-JUS/DGTAIPD exige, pero su
redacción final debe ser revisada por el área legal de LAP antes de su
colocación física.

Esa revisión tiene un contenido preciso: contrastar los campos del modelo
uno por uno contra el texto de la Directiva y contra el D.S.
016-2024-JUS, para confirmar que no falta ningún elemento informativo
exigido ni sobra alguno que la norma no requiera. Este informe propone el
conjunto de campos, no acredita que sea el exigible.

\subsection{Elementos que el modelo incorpora}

\begin{center}
\small
\begin{tabularx}{\linewidth}{@{}L{4.2cm}Y@{}}
\toprule
\textbf{Elemento} & \textbf{Razón de su inclusión} \\
\midrule
Identidad del responsable &
El titular debe saber ante quién ejercer sus derechos. \\
\addlinespace
Finalidad del tratamiento &
Núcleo del deber de información y límite frente al deslizamiento de
finalidad (Sec.~\ref{sec:impacto-etico}). \\
\addlinespace
Base de legitimación &
Permite al titular entender por qué no se le pidió consentimiento. \\
\addlinespace
Plazo de conservación &
Exigencia del principio de conservación limitada. \\
\addlinespace
Destinatarios &
Incluye la eventual entrega a autoridades conforme a las normas de
seguridad ciudadana (Sec.~\ref{sec:dl1218}). \\
\addlinespace
Canal para ejercer derechos &
Sin un canal indicado, el derecho es nominal. \\
\addlinespace
Declaración de lo que el sistema \emph{no} hace &
No es un elemento exigido por la norma. Se incorpora por decisión de
este frente: es la traducción visible del compromiso contra el
deslizamiento de finalidad y el elemento que distingue este cartel de
uno genérico de CCTV. \\
\bottomrule
\end{tabularx}
\end{center}

\subsection{Modelo propuesto}

En el modelo que sigue, los textos en \emph{cursiva} señalan los campos
que este informe no puede completar porque corresponden a datos que solo
LAP posee: su identificación registral, la base de legitimación que
decida invocar, sus canales de atención y el plazo de conservación que
fije. El resto del contenido es propuesta de este frente.

\vspace{0.5em}
\noindent
\fbox{%
\begin{minipage}{\dimexpr\linewidth-2\fboxsep-2\fboxrule\relax}
% Interlineado sencillo dentro del cartel: con el 1,5 del documento a
% 12 pt los filetes quedaban separados del texto por huecos muy grandes.
\setstretch{1}
\vspace{0.6em}
\begin{center}
{\Large\bfseries ZONA VIDEOVIGILADA}\\[0.25em]
{\large\bfseries SISTEMA DE ESTIMACIÓN DE AFORO}\\[0.2em]
{\small\itshape Este sistema NO identifica personas}
\end{center}
\noindent\rule{\linewidth}{0.8pt}
\vspace{0.2em}

\small
\setlength{\extrarowheight}{3pt}
\begin{tabularx}{\linewidth}{@{}L{3.1cm}Y@{}}
\textbf{Responsable} &
Lima Airport Partners S.R.L. (\emph{razón social registral completa,
RUC y domicilio fiscal, a completar por LAP}). \\

\textbf{Finalidad} &
Estimar el número de personas presentes por zona, con el fin de prevenir
aglomeraciones y gestionar el flujo de pasajeros en el terminal. \\

\textbf{Base legal} &
\emph{Base de legitimación que LAP invoque, con su fundamento normativo
exacto, a completar por su área legal.} \\

\textbf{Conservación} &
Las imágenes no se almacenan. El sistema conserva únicamente conteos
agregados por zona, durante el plazo fijado en su política de retención
(\emph{plazo a definir con LAP, medida M-05}). \\

\textbf{Destinatarios} &
Personal autorizado de LAP. Las imágenes del sistema de seguridad podrán
entregarse a las autoridades competentes cuando lo requieran conforme a
ley. \\

\textbf{Sus derechos} &
Puede ejercer sus derechos de acceso, rectificación, cancelación y
oposición escribiendo al canal de atención de LAP o acudiendo a su
oficina de atención al usuario (\emph{dirección de correo y ubicación
de la oficina, a completar por LAP}). \\

\textbf{Más información} &
\emph{Dirección web de la política de privacidad de LAP, a completar.} \\
\end{tabularx}

\vspace{0.4em}
\noindent\rule{\linewidth}{0.8pt}
\vspace{0.2em}

{\footnotesize\bfseries Este sistema no realiza reconocimiento facial, no
identifica a personas concretas, no realiza seguimiento individual de
trayectorias y no adopta decisiones automatizadas sobre los pasajeros.}

\vspace{0.6em}
\end{minipage}%
}
\vspace{0.8em}

\subsection{Criterios de colocación}

\begin{enumerate}[leftmargin=1.6em,itemsep=0.35em]
  \item \textbf{Anterioridad.} El cartel debe ser visible \emph{antes} de
        que la persona ingrese a la zona videovigilada, no dentro de
        ella. Informar a quien ya fue captado no cumple la finalidad de
        la norma.
  \item \textbf{Cobertura de todos los accesos.} Cada acceso a una zona
        cubierta requiere señalización. Un acceso sin cartel es un punto
        de incumplimiento con independencia de cuántos carteles haya en
        el resto del terminal.
  \item \textbf{Legibilidad.} Tamaño y contraste suficientes para ser
        leídos a la distancia de aproximación real, considerando que el
        pasajero circula con equipaje y no se detiene.
  \item \textbf{Idiomas.} Español e inglés como mínimo, dada la condición
        de terminal internacional. Si conviene añadir algún idioma más es
        algo que depende de la composición real del tráfico del
        aeropuerto, dato que debe pedirse a LAP y que este informe no
        supone.
  \item \textbf{Accesibilidad.} Altura y ubicación que permitan la
        lectura desde una silla de ruedas. El criterio se enuncia aquí en
        términos funcionales y no normativos: la normativa peruana de
        accesibilidad aplicable a señalización en establecimientos de uso
        público no se ha contrastado, de modo que las medidas concretas
        deben tomarse de ella y no de este anexo.
  \item \textbf{Registro de la colocación.} Documentar fotográficamente
        cada cartel instalado, con fecha y ubicación. Ante una
        fiscalización, el cumplimiento del deber de información debe
        poder acreditarse, y la fotografía fechada es la evidencia más
        simple de producir.
\end{enumerate}

\subsection{Observación final sobre este anexo}

Este cartel es el punto donde el análisis normativo de este informe se
vuelve visible para el pasajero. Todo lo demás (la arquitectura, las
medidas de privacidad desde el diseño, la auditoría) ocurre en
documentos que el titular de los datos nunca verá. La frase
\enquote{este sistema no identifica personas} solo es defendible si las
medidas M-02, M-03 y M-04 de la Sección~\ref{sec:privacidad-diseno}
están efectivamente implementadas.

Colocar un cartel que declare más de lo que el
sistema cumple es peor que no colocarlo: convierte una deficiencia
técnica en una declaración inexacta hacia el titular. La redacción final
del cartel debe fijarse \emph{después} de que la arquitectura esté
congelada, y debe revisarse si la arquitectura cambia.
